\chapter{数据清洗：让AI帮你，但别让AI替你决策}

\section{数据清洗的哲学：清洗决策是领域判断}

AI 可以执行填充、删除和类型转换等技术操作，但是否填充、删除或保留某条数据，需要结合领域知识判断。

\begin{enumerate}
  \item 速度为 0，可能是异常，也可能是信号交叉口、收费站、拥堵路段或事故现场的真实状态。
  \item GPS 轨迹跳点可能是定位错误，也可能是快速路切换。
  \item 网约车异常订单可能来自数据错误，也可能由平台策略造成。
\end{enumerate}

\section{缺失值处理：用Prompt描述业务逻辑}

让 AI 自动填充所有缺失值可能掩盖重要事件。例如，使用均值填充检测器故障期间的数据，会掩盖流量中断。

正确做法是先定义规则，再让 AI 实现。例如：

\begin{quote}
对于字段 speed，当 station status 为 fault 时，缺失值应标记为 NaN 并保留，不进行填充；当 station status 为 normal 且缺失率低于 5\% 时，可以使用前后 15 分钟均值插值。
\end{quote}

还应结合交通场景理解完全随机缺失、随机缺失和非随机缺失等不同缺失机制。

\section{异常值检测：与AI协作设计检测规则}

单纯使用统计规则可能误删真实的重要事件。例如，$3\sigma$ 原则可能删除大雨天的真实需求峰值。更可靠的策略是联合使用统计检测、业务规则校验和人工抽查。

\section{特征工程：将领域知识编码为特征}

\begin{itemize}
  \item 时间特征：工作日、节假日和高峰期。
  \item 空间特征：路段等级、周边兴趣点和路网拓扑。
\end{itemize}

特征工程的创意来自领域知识，具体实现可以交给 AI 辅助完成。

\section{数据标准化与编码：让AI处理，但要检查输出逻辑}

对小时数直接进行最小最大标准化会破坏时间的周期性，此时更适合使用正弦和余弦编码。类别特征应根据语义选择独热编码、标签编码或目标编码。

\section{实战练习：使用完整Prompt链清洗数据}

选择一份真实数据，先定义缺失值、异常值和编码规则，再让 AI 按照规则生成清洗代码，并逐步检查其输出逻辑。
